\scanpage{95}

% L1
\uline{Lemme} Soit $\unsure{\mathcal E|E}$ fibré vectoriel de rg $N\geq3$. $X=\check{\mathbb P}(\mathcal E)$.

% L2–3
a) Soit $\unsure{\tau_X|\mathrm T_X}$ une \inserted{ps.}réflexion de $X$. Alors $\exists1$ \inserted{ps.}réflexion $\tau_{\mathcal E}$ de $\mathcal E$ qui lui donne naissance, propre \unsure{ssi|\deleted{ssi}|si} $\tau_{\deleted{\unsure{X}}\mathcal E}$ l'est.

% L4–7
b) Soit $C$ un \deleted{hyperplan} \inserted{sous-fibré proj} de $X$, \deleted{stable pour $\tau_X$, de} \inserted{$C=\check{\mathbb P}(V)$, supposons} \deleted{$\tau_X C=C$ i.e. $\tau_{\mathcal E}V=V$.} $\tau_X$ propre, centre $h\in C$, \deleted{alors} alors \deleted{Considérons l'extension $\mathcal E$ de $\mathcal O_S$ par $V$} $\tau_{\mathcal E}$ induit l'identité sur $\mathcal E/V$.

% L7–11
Si $C$ est hyperplan et $\mathcal E$ l'enveloppe vectorielle can. de $E=X\smallsetminus C$, donc $\mathcal E/V\simeq\mathcal O_S$ can$^{\mathrm t}$. alors l'action de $\tau_X\mid E$ \deleted{sur $E$} s'identifie à $\tau_{\mathcal E}\mid E$, si on identifie $E$ à l'image inv. de $1$ par $\mathcal E\longrightarrow\mathcal O_S$

% L12–13
Dégager : condition de stabilité d'un ss-fibré projectif par pseudo réflexion

% L14–15
conditions de commutation de deux pseudo-réflexions.

% L16
Ordre d'une pseudo-réflexion

% L17
Ordre du produit de deux pseudo-réflexions

% archive sheet number: 94
